\section{Related Work}
\label{sec:related-work}

\paragraph{Feed-Forward 3D and 4D Reconstruction.}
Feed-forward approaches have demonstrated significant potential in directly processing multi-view images to infer 3D geometric attributes. Foundational architectures such as DUSt3R~\cite{wang2024dust3r} and VGGT~\cite{wang2025vggt} bypass traditional optimization by directly predicting 3D point clouds and depth maps. To further enhance efficiency, FastVGGT~\cite{shen2025fastvggt} introduces region-based random sampling, which markedly reduces inference latency, particularly for large-scale inputs. Concurrently, increasingly powerful models like Concurrently, models like Depth Anything V3~\cite{lin2025depth} and SAM3D~\cite{chen2025sam} have advanced generalizable 3D scene understanding. 
Beyond point clouds, MVSplat~\cite{chen2024mvsplat} and AnySplat~\cite{jiang2025anysplat} extend this paradigm to Gaussian Splatting, though primarily for static environments.
Recent works have begun to address temporal dynamics within this paradigm. StreamVGGT~\cite{zhuo2025streaming} enables streaming reconstruction by maintaining a memory cache of previously computed tokens and leveraging temporal causal attention to accelerate inference over time. In the driving domain, DriveVGGT~\cite{jia2025drivevggt} utilizes the VGGT backbone within a scale-aware framework to output depth and point clouds specifically for autonomous driving data.
Our work shifts the focus from static geometric design to enhancing dynamic scene representation and novel-view fidelity. By adapting the VGGT formulation specifically for the complexities of autonomous driving, we achieve significant gains in visual quality and temporal consistency through improved feature utilization.


\paragraph{Gaussian Splatting for Autonomous Driving Scene Reconstruction.}
Existing methods primarily rely on per-scene optimization, iteratively refining Gaussian kernels to minimize rendering loss. AutoSplat~\cite{khan2025autosplat} achieves high-fidelity reconstruction by imposing geometric constraints on varying regions, while PVG~\cite{chen2023periodic} introduces periodic vibration Gaussians for dynamic urban environments. To handle the complexity of traffic, StreetGaussians~\cite{yan2024street} decomposes scenes into static backgrounds and moving vehicles, and DrivingGaussian~\cite{zhou2024drivinggaussian} introduces a dynamic Gaussian graph for moving objects. Similarly, S3Gaussian~\cite{huang2024textit} utilizes a spatial-temporal field network to model 4D dynamics. Despite their high quality, these optimization-based methods require hours of computation per scene, making them unscalable for large-scale autonomous driving datasets.
In contrast, the feed-forward paradigm generalizes to diverse scenes without requiring per-scene optimization. Drivingforward~\cite{tian2025drivingforward} reconstructs adjacent frames through interpolation, while STORM~\cite{yang2024storm} employs a transformer to generate per-frame Gaussians and scene flow via learnable motion tokens. However, these methods do not leverage 3D foundation models, which limits their representative capacity. More recently, WorldSplat~\cite{zhu2025worldsplat} introduced a 4D-aware latent diffusion model for pixel-aligned Gaussian production, and DGGT~\cite{chen2025dggt} proposed diffusion-based refinement to mitigate motion artifacts. Nevertheless, the integration of diffusion models significantly increases inference latency.
Unlike these approaches, ReconDrive leverages a 3D foundation model to enable scalable 4D Gaussian prediction with the well-designed architecture. By incorporating explicit motion and temporal consistency, we achieve efficient urban scene reconstruction without the need for iterative optimization or high-latency post-refinement.